\documentclass[11pt]{ctexart}
\usepackage{amsmath,amssymb}
\begin{document}
极大无和集计数：完整证明草稿

\(f_{\max}(\mathbb{F}_5^k)=3^{5^k/10+o(5^k)}\)

研究证明档案 · 2026-11-22 · 纸面论证与精确有限证书

\section{结果、定义与稿件状态}

设 \(G=\mathbb{F}_5^k\)，\(n=5^k\)。称 \(A\subseteq G\) 为无和集，若不存在 \(x,y,z\in A\) 满足 \(x+y=z\)；这里允许变量重复。无和集的“极大”按包含关系定义。记 \(f_{\max}(G)\) 为极大无和集的数目。

\textbf{主定理（本稿的证明结论）。} 当 \(k\to\infty\) 时，
\[
 f_{\max}(\mathbb{F}_5^k)=3^{n/10+o(n)}.
\]
等价地，\(n^{-1}\log_3 f_{\max}(\mathbb{F}_5^k)\to 1/10\)。这正是 Hassler--Treglown 的 Conjecture 2.3 [1]。这里证明的是指数主项；未声称确定多项式前因子或全部极大集合的结构。

本稿由完整的全局归约、穷尽的局部分型和所需引理证明组成。正文后的技术附录收录这些引理的完整证明及有限证书表；源文件包同时提供重建图、转移矩阵和精确有理数不等式的程序。证明中不再保留未经证明的局部猜想作为假设。

\textbf{验证状态。} 配套 Lean 工程已形式化验证本稿的主定理。部分局部步骤另附整数或有理数证书，18 项精确有限检查均已通过。各技术附录保留成稿时的审阅记录。本稿现为待外部同行审阅的公开修订稿，文献优先权与新颖性仍待专门核查。

主要输入的既有文献是：Green--Ruzsa 容器定理的 [2, Lemma 2.3] 表述，Green 的三集合移除引理 [3, Theorem 1.5]，以及 Lev 的大无和集结构定理 [4, Theorem 3]。其余用于完成局部估计的论证均在本证明档案中明确给出。

\section{链接图与通用块}

写 \(G=\mathbb{F}_5\times H\)，\(h=|H|=n/5\)，并令 \(B=\{2,3\}\times H\)。对于无和集 \(S\subseteq G\setminus B\)，链接图 \(L_S[B]\) 的顶点集为 \(B\)。两个不同顶点 \(x,y\) 相邻，当且仅当 \(x+y\in S\)、\(x-y\in S\)、\(y-x\in S\) 至少一个成立；自环顶点集合为
\[
 \Lambda=B\cap\bigl((S+S)\cup(S-S)\cup S/2\bigr).
\]
因而 \(I\subseteq B\) 在该图中独立，当且仅当 \(S\cup I\) 无和。用 \(\operatorname{mis}(F)\) 表示图 \(F\) 的极大独立集数目。

记三个外部层的尾坐标为 \(A=S_0\)、\(P=S_1\)、\(Q=-S_4\)，并记
\[
 z=|A|,\quad p=|P|,\quad q=|Q|,\quad
 a_0=|A\cup(-A)|,\quad d_0=|P\cup Q|.
\]
为避免混淆，子空间的阶一律记作 \(\kappa\)，不记作 \(q\)。取包含所有尾坐标的子空间 \(K\le H\)，\(\kappa=|K|\)。对 \(g\notin K\)，通用块位于 \(\{2,3\}\times((g+K)\cup(-g+K))\)，共有 \(4\kappa\) 个顶点，没有自环。

以 \(X_x=(2,x)\)、\(Y_x=(3,-x)\) 标记其顶点。完整邻接条件是
\[
 X_xX_y:\quad x-y\in A\cup(-A)\ \text{或}\ x+y\in-Q;
\]
\[
 Y_xY_y:\quad x-y\in A\cup(-A)\ \text{或}\ x+y\in-P;
\]
\[
 X_xY_y:\quad x-y\in A\ \text{或}\ x+y\in-(P\cup Q).
\]
差条件保持陪集符号，和条件反转陪集符号，所以在通用块内没有这两类边的重合。恰好一半顶点的度数为 \(a_0+z+d_0+p\)，另一半为 \(a_0+z+d_0+q\)。

\textbf{通用块定理（GB）。} 对任意无和种子 \(S\subseteq\{0,1,4\}\times K\)，其通用块 \(\Gamma_S\) 满足
\[
 \operatorname{mis}(\Gamma_S)\le 3^\kappa.
\]
这个定理没有种子大小或仿射秩限制。第 3--5 节说明证明的主要机制，第 6 节及附录给出全部类型的覆盖。

\subsection{两个基本计数事实}

若 \(J\) 是简单图 \(F\) 的诱导子图，则 \(\operatorname{mis}(J)\le\operatorname{mis}(F)\)。证明：将 \(J\) 的每个极大独立集确定性地扩充为 \(F\) 的极大独立集；扩充集与 \(V(J)\) 的交仍是原集，故扩充映射是单射。这也说明，先删掉自环顶点再计数，不会超过保留全部顶点的普通边图的极大独立集数。

每个 \(N\) 点简单图满足 Moon--Moser 界 \(\operatorname{mis}(F)\le3^{N/3}\)。简短证明如下：取最小度数 \(\delta\) 的顶点 \(v\)，按闭邻域中第一个被选顶点分支。每支至少删去 \(\delta+1\) 个顶点，所以归纳给出 \((\delta+1)3^{(N-\delta-1)/3}\)。对整数 \(t\ge1\)，有 \(t\le3^{t/3}\)，从而结论成立。

\section{最后三类的熵证明}

在 \(A=\varnothing\) 时，通用块是二部图 \(\Gamma(P,Q)\)。左右两侧各有两个 \(K\) 索引的顶点排，分别记为 \(A_x,B_x\) 和 \(C_y,D_y\)。此处顶点符号 \(A_x\) 与上一节的集合 \(A\) 无关。令 \(T=P\cup Q\)，邻接为
\[
 A_xC_y:\ y-x\in T,\quad A_xD_y:\ y-x\in Q,
\]
\[
 B_xC_y:\ y-x\in P,\quad B_xD_y:\ y-x\in T.
\]
均匀随机选取一个极大独立集 \(I\)，定义四状态变量
\[
 Z_x=\mathbf{1}_{A_x\in I}+2\mathbf{1}_{B_x\in I},\qquad
 W_y=\mathbf{1}_{D_y\in I}+2\mathbf{1}_{C_y\in I}.
\]
二部图中的极大独立集由任一侧的选择唯一决定，故 \(H((Z_x)_{x\in K})=\log_2\operatorname{mis}\Gamma\)。平移以及 \(A_x\leftrightarrow D_{-x}\)、\(B_x\leftrightarrow C_{-x}\) 都是图自同构，因此全部 \(Z_x,W_y\) 同边际。这里不假设任何独立性。

\subsection{加权局部熵引理}

给定互异的 \(d\) 个平移，令 \(\boldsymbol z\) 为相应输入，\(F(\boldsymbol z)=W_y\) 为由独立性与极大性唯一确定的输出。若正权函数 \(R\) 满足 \(\mathbb E\log_2 R=0\)，且 \(\Omega\) 包含局部输入的全部支撑，令
\[
 \mathcal Z=\sum_{\boldsymbol z\in\Omega}R(\boldsymbol z).
\]
Jensen 不等式给出 \(H(\boldsymbol z)\le\log_2\mathcal Z\)：对正概率输入将 \(R(\boldsymbol z)/\Pr(\boldsymbol z)\) 代入凹函数 \(\log_2\) 即得。将全部平移后的输入窗口求和，每个坐标恰覆盖 \(d\) 次。链式法则和“增加条件不会增加熵”给出 Shearer 不等式，因而
\[
 \operatorname{mis}\Gamma\le\mathcal Z^{\kappa/d}.
\]
一种保证期望抵消的选权方式是
\[
 R=\prod_i\frac{w_i(z_i)}{w_i(F(\boldsymbol z))},
\]
因为每个输入与输出同边际。

输入状态写作 \(s=x+2y\)。一个 \(P\setminus Q\) 输入对 \((C,D)\) 的阻挡状态贡献 \((x\vee y,y)\)，一个 \(Q\setminus P\) 输入贡献 \((x,x\vee y)\)，公共输入贡献 \((x\vee y,x\vee y)\)。对所有输入逐坐标取或后得到 \((b_C,b_D)\)，输出就是 \((1-b_D)+2(1-b_C)\)。

\subsection{类型 A 与 C：唯一私有邻居}

若 \(Q\setminus P=\{t\}\)，则 \(N(A_x)\setminus N(B_x)=\{C_{x+t}\}\)。若 \(Z_x=2\)，极大性强迫 \(C_{x+t}\) 被选；反向私有邻居论证又给出逆命题。因此逐个极大独立集都有
\[
 Z_x=2\quad\Longleftrightarrow\quad W_{x+t}=2.
\]
在类型 A、C 中，其余三个平移都属于 \(P\)。将唯一 \(Q\)-专属输入放在最后，以上关系说明四输入支撑包含于
\[
 \Omega=\{\boldsymbol z\in\{0,1,2,3\}^4:
 z_4\ne2\ \text{或}\ (z_1,z_2,z_3)=(0,0,0)\},
\]
该集合恰有 193 个元素。

类型 A 满足 \(|P|=3,|Q|=1,P\cap Q=\varnothing\)。三个 \(P\) 输入取权 \((1,1/2,1/2,1/2)\)，\(Q\) 输入取权 \((1,1,1,1/2)\)。按输出状态 \(0,1,2,3\) 求和，四项依次为 \(571/16,19,8,16\)，所以
\[
 \mathcal Z_A=\frac{1259}{16}<81,
 \qquad\operatorname{mis}\Gamma\le
 \left(\frac{1259}{16}\right)^{\kappa/4}<3^\kappa.
\]
类型 C 满足 \(|P|=3,|Q|=2,|P\cap Q|=1\)。公共输入权为 \((1,1,1,1)\)，两个 \(P\)-专属输入权仍为 \((1,1/2,1/2,1/2)\)，\(Q\)-专属输入权为 \((1,1/2,1,1/2)\)。四项依次为 \(191/4,10,4,8\)，得到
\[
 \mathcal Z_C=\frac{279}{4}<81,
 \qquad\operatorname{mis}\Gamma\le
 \left(\frac{279}{4}\right)^{\kappa/4}<3^\kappa.
\]
两种论证只要求四个平移互异，不要求它们仿射独立。完整逐项求和见附录 E1。

\subsection{类型 B：共享输入的两个窗口}

现在 \(P=\{0,u\}\)、\(Q=\{v,w\}\) 不交。先设七个坐标
\[
 0,-u,u,-v,-w,u-v,u-w
\]
互异。以此顺序记七输入为 \(z_0,\ldots,z_6\)。两个右侧输出位于 \(0,u\)，共享输入 \(z_0\)，各另有一个 \(P\) 输入和两个 \(Q\) 输入，记为 \(f_0,f_1\)。使用正权
\[
 c=(1,1/4,1,1),\quad a=(1,1/3,2/3,1/2),
\]
\[
 b=(1,2/3,1/3,1/2),\quad g=(1,2/27,2/27,1/8).
\]
它们逐状态满足 \(ca^2b^4=g^2\)。取
\[
 R(\boldsymbol z)=
 \frac{c(z_0)a(z_1)a(z_2)b(z_3)b(z_4)b(z_5)b(z_6)}
 {g(f_0)g(f_1)}.
\]
七输入和两输出同边际，故 \(\mathbb E\log_2R=0\)。对共享状态条件化后，完整的 \(4^7\) 项求和分解成四项：
\[
 R_0=\frac{2629}{72},\quad R_1=\frac{775}{24},\quad
 R_2=R_3=\frac{125}{8}.
\]
这里每个 \(R_s\) 是单个三叶窗口对输出权取商后的加权和。其直接计算只需原始叶权总质量 \((5/2)^3=125/8\)，以及当 \(s=0\) 时输出 \(3,1,2\) 的原始质量分别为 \(1,1/3,7/9\)，当 \(s=1\) 时输出 1 的原始质量为 \(4/3\)。其余输入输出皆为 0。
\[
 \mathcal Z_B=\sum_{s=0}^3c(s)R_s^2
 =\frac{43177189}{20736}<2187=3^7.
\]
七重覆盖遂给出 \(\operatorname{mis}\Gamma\le\mathcal Z_B^{\kappa/7}<3^\kappa\)。

\subsection{类型 B 的坐标碰撞与低秩}

七坐标发生碰撞，当且仅当 \(v\) 或 \(w\) 属于 \(\{-u,2u\}\)，或者 \(w-v\in\{u,-u\}\)。直接比较 \(\{0,-u,u\}\)、\(\{-v,u-v\}\)、\(\{-w,u-w\}\) 三个集合即可验证这一分类。

若后一条件成立，则 \(Q=P+t\)。记 \(C_R\) 为平移集合 \(R\) 的 Boolean 矩阵，则在 Boolean 乘法下有
\[
 M(R+U,R+V)=\operatorname{diag}(C_R,C_R)M(U,V).
\]
二部图的极大独立集与其邻接矩阵的不同“行并集”一一对应。乘积 \(XY\) 的任意行并集都是 \(Y\) 的行并集，所以 \(\operatorname{mis}\Gamma(P,P+t)\le\operatorname{mis}\Gamma(\{0\},\{t\})\)。这用的是行并集的包含关系，不是加边的单调性。

若只发生前一类碰撞，且仿射秩至少二，假设 \(v\in\{-u,2u\}\) 而 \(w\notin\langle u\rangle\)。交换 \(P,Q\) 再平移 \(-v\)，得到 \(P'=\{0,w-v\}\)、\(Q'=\{-v,u-v\}\)。新 \(P'\) 方向在线外、\(Q'\) 在线上，两种碰撞条件都不成立。剩下的情形秩一。

秩一只需一个 20 点计数。对 \(K=\mathbb F_5\) 及非空 \(P,Q\)，任取 \(p\in P,q\in Q\)，五组顶点 \(A_x,B_x,C_{x+p},D_{x+q}\) 构成覆盖全图的五个 \(K_{2,2}\)。若某个极大独立集占据 \(a\) 个左配对、\(b\) 个右配对，则 \(a+b\le5\)。至少一边占据至多两个配对，而单边决定全部极大独立集，故
\[
 \operatorname{mis}\Gamma\le2\sum_{j=0}^{2}\binom{5}{j}3^j
 =212<243=3^5.
\]
平行情形化得的两平移图按五阶方向分成此类块；秩一情况亦逐陪集使用该界。至此类型 B 在所有秩上闭合，不需要秩二的大型搜索证书。完整证明见附录 E2。

\section{固定势函数与含零层的归纳}

令 \(c_0=\log_2(3)/3\)。对简单图 \(F\) 定义
\[
 \mu_\beta(F)=(c_0+\beta)|V(F)|-\beta|E(F)|.
\]
在整棵归纳树上固定 \(\beta\)，不随残余图重新选择。若当前最大度为 \(d\)，在该顶点作排除／选取分支，势函数下降量至少为
\[
 A(d)=c_0-\beta(d-1),\qquad
 B(d)=(d+1)c_0-\beta(d^2-d-1).
\]
第二式来自：与闭邻域相交的边数至多为 \(d^2\)。因此只需控制 \(2^{-A(d)}+2^{-B(d)}\)。这两个指数之和关于实数 \(d\) 为凸函数；度数 4 到固定上限 \(D\) 只需检查端点。三度以下的图由叶子、悬挂三角形、二三度相邻点、三正则图、环和独立分量分别处理。

\textbf{一般图主引理。} 若 \(0<\beta\le1/25\)、\(D\ge4\)，且上述端点和在 \(D\) 处至多为 1，则所有最大度至多 \(D\) 的图都有 \(\operatorname{mis}(F)\le2^{\mu_\beta(F)}\)。附录 P1 给出包括所有低度数情形的完整归纳和精确整数不等式。

由此得到两个用于本稿的结论：任何 \(d\ge10\) 的正则图有 \(\operatorname{mis}(F)<3^{|V(F)|/4}\)；若一半顶点度为 \(a\)，另一半为 \(b\)，且 \(a\ge b\ge10\)、\(a\le2b-3\)，则
\[
 \operatorname{mis}(F)\le3^{|V(F)|/4}.
\]
第二个结论取固定 \(\beta=c_0/(a+b-4)\)。最大度端点先对 \(a\) 作凸性约化，再对 \(b\) 用单调性，最终只需精确检查 \((a,b)=(10,10),(17,10)\)。初始势函数恰为 \(|V(F)|\log_2(3)/4\)。

若种子含零层且 \(|S|\ge8\)，则取两种度数 \(a\ge b\)，有 \(b\ge a_0+|S|\ge10\) 及 \(a\le2b-3\)，所以全部由该结论处理。剩下大小至多七的类型由附录分类处理。

\subsection{无三角形图的势函数}

两个补充估计将无零层种子进一步压缩到小类型。对无三角形图，令
\[
 \nu_b(F)=\frac{|V(F)|}{2}
 -b\left(|E(F)|-\frac{|V(F)|}{2}\right).
\]
当最大度至多 10 时，可取 \(b=1/36\)；当最大度至多 \(D\ge12\) 时，可取 \(b=2/[7(D-1)]\)，均有 \(\operatorname{mis}(F)\le2^{\nu_b(F)}\)。边界 \(D=12\) 的三层分支系数为 \(9997293/10000000<1\)，因此该边界没有被略去。完整证明见附录 P2。

另有对所有最大度至多 8 的无三角形图成立的估计
\[
 \operatorname{mis}(F)\le2^{(16|V(F)|-|E(F)|)/31}.
\]
其度数八的正则情况使用一次有序邻域展开，完整证明见附录 P3、P4。对最大度至多九、且每个九度顶点至少有四个度数至多七的邻居的无三角形图，同一势函数仍适用；这个类对诱导删点封闭，其证明见附录 Z4 第 2 节。

这些引理都先对所需的所有诱导子图证明，因而未假设删点之后仍保持原来的正则性或两种度数比例。对少数含三角形的低度类型，额外使用保留的原始邻域边、六个公共邻居和两个保护邻居；其逐步坐标核查及残余图归纳在附录 Z3--Z5。

\section{结构引理与有限转移证书}

对称种子 \(S=-S\) 使合法顶点按反向配对成为假孪生点；极大独立集必须整对选择，商图的 Moon--Moser 界即给出所需估计。若只增加一个未配对元素，则把占据的配对支撑编码成“支配且诱导红度至多一”的集合。

精确的支撑计数引理为：设蓝图 \(Q\) 任意、红图 \(R\) 无三角形；同时满足蓝独立、红诱导最大度至多一、支配 \(Q\cup R\) 的集合数至多 \(3^{|V|/2}\)。最小度数分支并以红度破除平局给出完整归纳，见附录 S2。两个未配对元素分别在外侧相反层时，两种红匹配交替，仍可用同一引理，见附录 S3。

较小的通用块由五边形、圆梯图、嵌套邻域及一个折叠不等式处理。例如 \(\operatorname{mis}(C_5\mathbin{\square}K_2)=10\)、\(\operatorname{mis}(C_{10}\mathbin{\square}K_2)=122\)。证明直接编码各梯级为左端、右端或空；空梯级必须由两侧相反状态支配。双零层独立方向的情况由五环折叠化成 50 点图的短计数。这些均有纸面证明，见附录 S1。

含零层的余下有限模型使用“选择状态、尚需下一片支配的状态”构成转移矩阵。将图分成循环切片；状态为 \((J,U)\)，其中 \(J\) 为本片所选的独立集，\(U\) 是尚未由本片及前片支配的顶点。转移要求相邻片选择相容、下一片支配 \(U\)，并精确更新下一片的未支配集合。对这些模型，循环闭步与极大独立集一一对应。

对一个周期转移矩阵 \(T\)，若有严格正整数向量 \(w\) 满足 \(Tw\le\lambda w\)，则对角相似变换后的每一行和至多为 \(\lambda\)，故 \(\operatorname{tr}(T^m)\le r\lambda^m\)，其中 \(r\) 是状态数。以下数值均为直接给出的精确证书，不是浮点特征值猜测。

同层一零两非零的 100 点模型：\(105\cdot8^{10}=112742891520<3^{25}\)。相反层模型：\(120\cdot8^{10}=128849018880<3^{25}\)。

大小四的三个 100 点模型依次有界 \(70\cdot40^5=7168000000\)、\(295\cdot65^5=342285734375\)、\(75\cdot36^5=4534963200\)，均小于 \(3^{25}=847288609443\)。大小五的剩余含三角形模型有界 \(185\cdot64^5=198642237440<3^{25}\)。

无零层三种子的仿射秩二模型另用行并输出的转移矩阵。其 83 个状态记录两层膨胀集合，转移条目计数不同的输出标签；每个全图行并输出都有带标签的闭步表示，因此极大独立集数由闭步数上界控制。该矩阵有 19 个旋转轨道；对角平衡后的 Frobenius 范数平方为 \(1433053/32<3^{10}\)，由矩阵范数界得闭步数小于 \(3^{25}\)。其秩一情形使用上述 \(212<243\) 纸面界。含零层有限模型的全部相依参数则降为明列的 20 点图，最深七层的包含／排除分支加 Moon--Moser 叶界，给出的最大上界为 \(236<243\)。

完整状态定义、切片与原链接图的逐边同构、正权表、20 点分支表均收入附录 F1--F5。含零层模型的验证器从有限域的原始和／差条件独立重建图；无零层 83 状态模型的验证器则分别用选行公式和显式最大原像判据构造全部转移条目，并逐项比对。

\section{全部种子的穷尽覆盖}

这一节的作用是确认没有在“通用情形”之外留下依赖关系或特殊秩。完整分类文件为附录 C0、C1、C2；下面概述其范围。

\textbf{不含零层。} 纯外侧层以及 \(P,Q\) 嵌套的情况任意大小均可直接处理。嵌套时每个左顶点配对的两个邻域有包含关系，极大性使左配对仅有三个可行状态，故至多 \(3^\kappa\)。大小至少七的其余种子由无三角形势函数处理。大小至多三由纸面小图界及 83 状态转移处理。

在大小四、五、六且 \(p\ge q>0\) 时，记 \(d=|P\cup Q|\)。全部分型如下：

大小四：\((p,q,d)=(3,1,3)\) 嵌套，\((3,1,4)\) 为熵类型 A；\((2,2,2)\) 嵌套，\((2,2,3)\) 为相反层双缺陷，\((2,2,4)\) 为熵类型 B。

大小五：\((4,1,4)\) 嵌套，\((4,1,5)\) 用九度低邻点引理；\((3,2,3)\) 嵌套，\((3,2,4)\) 为熵类型 C，\((3,2,5)\) 用八度势函数。

大小六：\((5,1,5)\) 嵌套，\((5,1,6)\) 用十二度势函数；\((4,2,4)\) 嵌套，\((4,2,5)\) 用九度低邻点引理，\((4,2,6)\) 用十度势函数；\((3,3,3)\) 嵌套，\((3,3,4)\) 用相反层双缺陷，\((3,3,5)\) 用八度势函数，\((3,3,6)\) 用十度势函数。这穷尽 \(p+q\) 和可能交集大小。

\textbf{含零层。} 纯零层图为 \(d=|A\cup(-A)|+|A|\) 正则图。\(d\ge10\) 用一般势函数；\(d\le9\) 迫使 \(|A|\le4\)。除对称、单缺陷和双种子外，只剩三个互不相反的零层种子，图为九正则。若该图无三角形，附录 P2 的十度势函数给出 \(2^{7N/18}<3^{N/4}\)，其中 \(N\) 为图的顶点数，严格比较等价于 \(2^{14}<3^9\)。若含三角形，则必有 \(A=\{a,b,-a-b\}\)，其每个开放邻域含至少六条边，另一固定势函数给出所需界，见附录 Z1、Z2。

混合含零层种子按大小处理：至少八用两度数主引理；六、七由原始度数上限为 \(12/9\)、\(10/9\) 的遗传引理及无三角形引理覆盖；五由九正则公共邻居、\(10/8\) 度数亏损归纳和一个 100 点转移模型覆盖；四由穷尽表、三个 100 点模型和双缺陷覆盖；三由 \(8/7\) 无三角形图、两个 100 点模型及全部相依的 20 点图覆盖；至多二由低秩纸面计数覆盖。

在使用全链接图版本的结构引理时，可用放大消除通用块与核心的区分：将固定种子嵌入 \(K\oplus\mathbb F_5^m\)，通用块重复 \((5^m-1)/2\) 次。全图界 \(3^{\kappa5^m/2}\) 迫使单块计数至多 \(3^\kappa\)。因而每个全图结构引理都合法地产生 GB 所需的单块结论。

上述各类与第 3 节的三个全秩熵证明合并，即证明 GB。

\section{从通用块到统一链接图上界}

\textbf{大度数指纹引理。} 若简单图 \(F\) 有 \(N\) 个顶点、最小度至少 4096，且最大度至多最小度的两倍，则
\[
 \operatorname{mis}(F)<3^{\theta N},\qquad
 \theta=\frac{1495048}{6041601}<\frac14.
\]
证明如下。固定顶点次序，对极大独立集 \(I\) 从空集合 \(T\) 开始。令 \(A=V(F)\setminus N(T)\)，只要 \(I\) 中存在在 \(F[A]\) 内度数至少 512 的点，就按固定次序将一个这样的点加入 \(T\)。每次至少删去 512 个尚未删去的邻点，且 \(T\) 与其邻域不交，故 \(|T|\le N/513\)。

停止后令 \(U=\{v\in A:d_{F[A]}(v)\le511\}\)，则 \(I\subseteq U\)，且 \(I\) 是 \(F[U]\) 的极大独立集。\(T\) 唯一确定 \(U\)。记最小度、最大度为 \(\delta,\Delta\)，跨边计数给出
\[
 (\delta-511)|U|\le\Delta(N-|U|),\qquad
 |U|\le\frac{8192}{11777}N.
\]
候选指纹数至多 \(\sum_{j\le N/513}\binom Nj<3^{8N/513}\)，而每个 \(U\) 的极大独立集数至多 \(3^{8192N/35331}\)。二者相乘即得，因为 \(8192/35331+8/513=\theta\)。二项式熵估计及全部常数的初等核查见附录 G。

对完整链接图的普通边图，先只看差边：它是 \(d=a_0+d_0\) 正则图。额外和边在任一顶点至多增加 \(z+\max(p,q)\le d\) 个邻点，所以
\[
 \delta\ge d\ge|S|/2,\qquad \Delta\le2d\le2\delta.
\]
因此 \(|S|\ge8192\) 时，删掉自环顶点的诱导图单射与指纹引理给出 \(\operatorname{mis}(L_S[B])<3^{n/10}\)。

若 \(|S|\le8191\)，取全部尾坐标张成的 \(K\)，则 \(\kappa\le5^{8191}\)。差边保持剩余陪集坐标，和边将其取负；全部自环位于零陪集。链接图遂分解成一个至多 \(2\kappa\) 点的核心，以及 \((h/\kappa-1)/2\) 个通用块。Moon--Moser 与 GB 给出
\[
 \operatorname{mis}(L_S[B])
 \le3^{2\kappa/3}(3^\kappa)^{(h/\kappa-1)/2}
 =3^{h/2+\kappa/6}.
\]
令绝对常数 \(C=5^{8191}/6\)，合并两种大小区间即得对所有种子统一成立的估计
\[
 \operatorname{mis}(L_S[B])\le3^{n/10+C}.
\]
常数虽不优化，但与 \(k\)、种子以及容器均无关。这一统一性正是下一节所需的条件。

\section{容器归约及主定理的上界}

由 [2, Lemma 2.3]，存在 \(2^{o(n)}\) 个容器覆盖所有无和集，每个容器仅有 \(o(n^2)\) 个 Schur 三元组，且两个误差都统一。由 [3, Theorem 1.5] 对三个集合 \(F,F,-F\) 应用移除引理，可从每个容器删除统一的 \(o(n)\) 个元素，使余下部分无和。

将这个无和部分扩充为全群中极大的无和集 \(M\)，并将例外集改成 \(E\setminus M\)，得到覆盖全部无和集的至多 \(2^{o(n)}\) 对 \((M,E)\)，其中 \(M\cap E=\varnothing\)、\(|E|=o(n)\)。固定一对，任一全局极大无和集 \(X\subseteq M\cup E\) 写成 \(S=X\cap E\)、\(I=X\cap M\)。全局极大性说明 \(I\) 是 \(L_S[M]\) 的极大独立集。

若 \(|M|\le3n/10\)，Moon--Moser 立即给出 \(\operatorname{mis}(L_S[M])\le3^{n/10}\)。若 \(|M|>3n/10\)，Lev 的结构定理 [4, Theorem 3] 给出
\[
 M\subseteq(e+H_0)\cup(-e+H_0)
\]
其中 \(H_0\) 为真子空间、\(e\notin H_0\)。大小迫使其余维为一；右侧本身无和，所以 \(M\) 的全局极大性迫使等号。通过线性同构可将 \(M\) 化为 \(\{2,3\}\times\mathbb F_5^{k-1}\)，于是第 7 节的统一界适用。

每个 \(E\) 仅有 \(2^{o(n)}\) 个候选种子，因而
\[
 f_{\max}(G)\le2^{o(n)}2^{o(n)}3^{n/10+C}
 =3^{n/10+o(n)}.
\]
更具体地，若容器的统一三元组密度是 \(t_n=(\log n)^{-1/9}\)，移除误差函数为 \(\rho(t)\to0\)，则上界指数可写成
\[
 \frac n{10}+C+\log_3(2)\,n
 \bigl((\log n)^{-1/18}+\rho(t_n)\bigr).
\]
该式显式说明所有误差对容器和种子一致。

\section{匹配下界}

取 \(S=\{(1,0)\}\)。链接图中 \((2,0),(3,0)\) 带自环；其余每个非零反向尾坐标对 \(\{x,-x\}\) 恰给出一条四点路
\[
 (2,x),\ (3,x),\ (3,-x),\ (2,-x).
\]
不同尾坐标对之间无边。每条四点路有三个极大独立集，所以
\[
 \operatorname{mis}(L_S[B])=3^{(h-1)/2}=3^{n/10-1/2}.
\]
对其中每个极大独立集 \(I\)，将 \(S\cup I\) 扩充为全群的极大无和集 \(X_I\)。扩充不能加入其他 \(B\) 顶点，故 \(X_I\cap B=I\)，这给出单射。因此 \(f_{\max}(G)\ge3^{n/10-1/2}\)。结合上界，主定理得证。这个下界也是 [1] 所陈述的构造，本稿将其写成链接图中的明确单射。

\section{复核入口与证明依赖}

附录按下列顺序组织：C 为穷尽分类，S 为结构及低秩纸面计数，P 为一般图和无三角形图势函数，Z 为含零层的余下分类与特殊归纳，F 为有限转移和 20 点分支证书，E 为最后三类全秩熵证明，G 为统一全局拼接。

源包保留各证明文件、精确证书、标准库验证器和统一复核入口 \texttt{verify\_all.py}。复核入口运行的是所列有限断言；它不把数学文本自动转成形式证明。正文和附录中的归纳、双射、分类穷尽及文献应用仍是必须人工审阅的数学内容。

较早探索中使用的 29 份秩二搜索 DAG 不属于最终证明依赖，未用来代替类型 B 的退化情况论证。证明也未使用“增加普通边不会增加极大独立集数”的错误命题，或把随种子变化的误差常数直接写成统一的 \(o(n)\)。

\section{参考文献}

[1] N. Hassler and A. Treglown, \emph{Notes on sum-free sets in abelian groups}, Integers 26 (2026), A97. Conjecture 2.3 and its lower-bound construction. 可核对作者版本 https://arxiv.org/html/2506.17401v2 。

[2] N. Hassler and A. Treglown, \emph{On maximal sum-free sets in abelian groups}, Electronic Journal of Combinatorics 29(2) (2022), P2.32. Lemma 2.3（容器）、Lemma 2.8（链接图）、Lemma 2.10（提升分解）。作者全文 https://web.mat.bham.ac.uk/A.Treglown/EJCaccepted.pdf 。

[3] B. Green, \emph{A Szemeredi-type regularity lemma in abelian groups, with applications}, Geometric and Functional Analysis 15 (2005), 340--376. 使用 Theorem 1.5 的三集合版本，取 \(F,F,-F\)。https://arxiv.org/pdf/math/0310476 。

[4] V. F. Lev, \emph{Large sum-free sets in \(\mathbb Z_5^n\)}, Journal of Combinatorial Theory, Series A 205 (2024), 105865. DOI 10.1016/j.jcta.2024.105865. 使用作者稿 Theorem 3，阈值为 \(3n/10\)。https://math.haifa.ac.il/seva/Papers/SFF5n.pdf 。

\end{document}
